\documentclass[10pt,a4paper]{amsart}
\usepackage[margin=19mm]{geometry}
\usepackage{amsmath,amssymb,mathtools}
\usepackage[scr=rsfs,cal=euler]{mathalfa}
\usepackage{xfrac}
\usepackage[unicode,hidelinks,bookmarks=false]{hyperref}
\newcommand{\ie}{i.e.\ }
\newcommand{\cf}{cf.\ }
\newcommand{\resp}{respectively\ }
\newcommand{\Subsection}{\subsection}
\providecommand{\nobreakdash}{\nobreak\hbox{-}}
\setlength{\emergencystretch}{2em}
\multlinegap0pt
\usepackage{amssymb}
\usepackage{enumerate}
\newtheorem{lem}{Lemma}
\newcommand{\beas}{\begin{eqnarray*}}
\newcommand{\eeas}{\end{eqnarray*}}
\title{\LARGE U\Large\MakeLowercase {niqueness of entire functions sharing two values with their partial derivative operators}}
\author{Sujoy Majumder, Debabrata Pramanik, Shantanu Panja}
\date{Source: arXiv:2605.08748v1 (CC BY 4.0)}
\begin{document}
\maketitle

\noindent\emph{Source and licence.} \LARGE U\Large\MakeLowercase {niqueness of entire functions sharing two values with their partial derivative operators}. Version arXiv:2605.08748v1 (CC BY 4.0), distributed by arXiv under the Creative Commons Attribution 4.0 International licence, http://creativecommons.org/licenses/by/4.0/, as recorded in the arXiv licence metadata for this exact version and shown on the abstract page https://arxiv.org/abs/2605.08748v1. Full licence notice: \texttt{licenses/arxiv-2605.08748-CC-BY-4.0.txt}.\par\smallskip

\noindent\textit{Editorial scope.} Counted unit: the article's lem ln2.11 (source0.tex lines 550--551). The article's complete original proof of it is delivered with it (source0.tex lines 552--598, one \texttt{proof} environment of 47 lines). No mathematics is written, repaired or re-derived: the statement and the proof are the article's own lines, in the article's own notation, and no retained line is substituted or re-flowed.  Retained uncounted as the source-local context the proof needs: lines 352--360; lines 544--548.  Nothing was omitted from any retained span, so the package carries no omission record.  The retained lines cite 3 works, whose entries are transcribed verbatim from the article's own bibliography source in the e-print and delivered with short editorial labels recorded in the item's reference mapping.  No retained line contains a figure environment, an image-inclusion command or drawing code, so no image is delivered and the package declares no asset dependency.  Editorial formatting only, in addition: an AMS \texttt{amsart} preamble that restates the article's own declarations and macros used by the retained lines, namely the environments \texttt{lem} on separate counters, together with the standard packages those lines need.  The retained environments print in this document as Lemma 1 in order of appearance, whereas the article prints the same items with the numbering of its own full document.  The complete native source of the article as distributed in the arXiv e-print for this exact version is bundled unchanged as \texttt{sources/200173/source0.tex}.
\begin{lem}\label{ln2.6}\cite[Lemma 3.59]{Hu-Li-Yang-2003} Let $P$ be a non-constant entire function in $\mathbb{C}^n$. Then 
\begin{align*}
\rho(e^P)=
\begin{cases}
\deg(P), & \text{if $P$ is a polynomial,}\\
+\infty, & \text{otherwise.}
\end{cases}
\end{align*}
\end{lem}

\begin{lem}\label{ln2.10}\cite{Majumder-Sarkar-2026} Let $f$ be a transcendental meromorphic function in $\mathbb{C}^n$ of finite order $\rho$. 
Then there exists an $\varepsilon$-set $E$ in $\mathbb{C}^n$ such that
\beas \displaystyle \left|\frac{\partial_{z_i}(f(z))}{f(z)}\right|\leq ||z||^{\rho-1+\delta},\eeas
holds for all large values of $||z||$, where $z\in \mathbb{C}^n\backslash E$ and $\delta>0$ is a given constant. 
\end{lem}

\begin{lem}\label{ln2.11} Let $f$ be a non-constant entire function in $\mathbb{C}^n$ of finite order. For $a\in\mathbb{C}^n$, if $f=a\rightleftharpoons \partial_{z_i}(f)=a$, where $i\in\{1,2,\ldots,n\}$, then $\partial_{z_i}(f)-a=A_i(f-a)$ for some non-zero constant $A_i$ in $\mathbb{C}$.
\end{lem}

\begin{proof} Note that $f=a\rightleftharpoons \partial_{z_i}(f)=a$, where $f$ is a non-constant entire function in $\mathbb{C}^n$ of finite order $\rho$. Using Lemma \ref{ln2.6}, we deduce that
\begin{align}\label{ln11.0}
\frac{\partial_{z_i}(f)-a}{f-a}=e^{P_i},
\end{align}
where $P_i$ is a polynomial in $\mathbb{C}^n$. We want to prove that $P_i$ is a constant. Suppose on contrary that $P_i$ is non-constant and let $\deg(P_i)=m_i$. Set
\begin{align*}
P_i(z)=\sum\limits_{\alpha_1,\ldots,\alpha_n=0}^{m_i}a_{\alpha_1\ldots\alpha_n}z_1^{\alpha_1}\ldots z_n^{\alpha_n}.
\end{align*}

Define
\begin{align*}
||a_{|\alpha|}||_{1}=\sum\limits_{\alpha_1+\ldots+\alpha_n=|\alpha|} |a_{\alpha_1\ldots\alpha_n}|.
\end{align*}

Now one can easily prove that (see the proof of Satz 1.2 \cite{Jank-Volkmann-1985})
\begin{align}\label{ln11.0a}
(1-o(1))||a_{|\alpha|}||_1r^{m_i}\leq |P_i(z)|\leq (1+o(1))||a_{|\alpha|}||_1r^{m_i}
\end{align}
for large values of $r=\|z\|$, where $|\alpha|=m_i$.

If we take $g=f-a$, then from (\ref{ln11.0}), we obtain
\begin{align}\label{ln11.3}
\frac{\partial_{z_i}(g)}{g}-\frac{a}{g}=e^P.
\end{align}

Let $A=\operatorname{supp}\mu^0_g$. Clearly $A$ is either empty or an analytic set of pure dimension $n-1$. Then the Lebesgue measure of $A$ in $\mathbb{C}^n$ is zero. Now by Lemma \ref{ln2.10}, there exists an $\varepsilon$-set $E$ in $\mathbb{C}^n$ such that $A\subset E$ and
\begin{align}\label{ln11.1} 
\left|\frac{\partial_{z_i}(g(z))}{g(z)}\right|\leq ||z||^{\rho-1+\delta}
\end{align}
holds for all large values of $||z||$, where $z\in \mathbb{C}^n\backslash E$ and $\delta>0$ is a given constant.
Let $S=\{\|z\|=r: z\in E,\quad \|z\|>1\}$. Then $S$ is of finite logarithmic measure. Choose a sequence $\{z_k\}$ such that $\|z_k\|=r_k$, $r_k\not\in [0,1]\cup S$ and $|g(z_k)|=M(r_k,g)$, $r_k\to +\infty$ as $k\to +\infty$. Clearly
\begin{align}\label{ln11.2} 
\lim\limits_{k\rightarrow +\infty}\frac{|a|}{|g(z_k)|}=\lim\limits_{k\rightarrow +\infty}\frac{|a|}{M(r_k,g)}=0,\;\;\text{i.e.,}\;\;\frac{|a|}{|g(z_k)|}=o(1)
\end{align}
for sufficiently large $\|z_k\|=r_k\not\in [0,1]\cup S$. 

Taking the principal branch of the logarithm, we deduce from (\ref{ln11.3}) that
\begin{align}\label{ln11.4}
P(z)=\log \left( \frac{\partial_{z_i}(g(z))}{g(z)} - \frac{a}{g(z)} \right)
=\log \left| \frac{\partial_{z_i}(g(z))}{g(z)}-\frac{a}{g(z)}\right|+ i\arg \left(\frac{\partial_{z_i}(g(z))}{g(z)}-\frac{a}{g(z)}\right).
\end{align}

Substituting (\ref{ln11.0a}), (\ref{ln11.1}) and (\ref{ln11.2}) into (\ref{ln11.4}), we get
\begin{align*} ||a_n||_1 r_k^{m_i}(1 - o(1))\leq |P(z_k)|&\leq \log \left|\frac{\partial_{z_i}(g(z_k))}{g(z_k)}\right|+\log \left(\frac{|a|}{|g(z_k)|}+e\right)+O(1)\\&\leq (\rho-1+\delta)\log r_k+O(1),
\end{align*}
which is impossible. Hence $P_i$ is a constant. If we take $A_i=e^{P_i}$, then from (\ref{ln11.0}), we have $\partial_{z_i}(f)-a=A_i(f-a)$ for some non-zero constant $A_i$ in $\mathbb{C}$.
\end{proof}

\begin{thebibliography}{99}
\bibitem[Hu-Li-]{Hu-Li-Yang-2003} {\sc P. C. Hu}, {\sc P. Li} and {\sc C. C. Yang}, Unicity of Meromorphic Mappings. Springer, New York (2003).
\bibitem[Jank-V]{Jank-Volkmann-1985} {\sc G. Jank} and {\sc L. Volkmann}, Einfuhrung in die Theorie der ganzen und meromorphen Funktionen mit Anwendungen auf Differentialgleichungen, Birkhauser, Basel-Boston, 1985.
\bibitem[Majumd]{Majumder-Sarkar-2026} {\sc S. Majumder} and {\sc N. Sarkar}, Bergweiler-Langley lemmas in several complex variables, \textit{Bull. Belgian Math. Soc.} (accepted for publication).
\end{thebibliography}
\end{document}
